% ==============================================================================
% פרק א': הגדרת המחקר, מסגרת תיאורטית והטמעת מתודולוגיית הדוח המכין (~6 עמודים)
% ==============================================================================

\chapter{הגדרת המחקר, מסגרת תיאורטית והטמעת מתודולוגיית הדוח המכין}
\label{ch:framework}

\begin{managerbox}[למה זה חשוב למנהל ולקובע המדיניות]
כל החלטת אשראי, רגולציה או השקעה המתקבלת לפני שהבעיה הוגדרה כמותית היא הימור. על פי מתודולוגיית הדוח המכין (\en{Preliminary Design Report}) \cite{segal2026prep}, תפקידו של פרק הפתיחה הוא להחליף הערכות איכותניות עמומות בחוזה מחקרי מדויק: להגדיר במספרים את פער הסיכון בשוק הדיור הישראלי (\num{2022--2026}), לגזור ממנו שאלת מחקר והשערות ניתנות להפרכה (\en{$H_1, H_2, H_3$}), לעגן אותן במטריצת ספרות ובמאמר דגל שפיט, ולתרגם אותן למפרט דרישות כמותי (\en{REQ-RES}) שלפיו תיבחן הצלחת המחקר.
\end{managerbox}

\section{מבנה המשפך: מהמגמה המאקרו-כלכלית אל פרויקט המחקר}
\label{sec:ch01_funnel}

דוח מכין מקצועי נפתח בסקירה תמטית הבנויה כמשפך (\en{Funnel Structure}), המוביל את הקורא מהמגמה העולמית והמאקרו-כלכלית הרחבה, דרך ההקשר המוסדי והיישומי בשוק המקומי, אל המוטיבציה המחקרית הממוקדת ואל מבנה הדוח \cite{segal2026prep,ulrich2020product}. איור~\ref{fig:ch01_funnel} מציג ויזואלית את ארבע תחנות המשפך של מחקר זה, המגשרות בין זעזוע הריבית הגלובלי לבין אמידת הסיכון ב"קיר המסירות" (\en{The Delivery Wall}) של שוק הדיור בישראל בשנים \num{2026--2027}.

\begin{figure}[htbp]
\centering
\begin{english}
\begin{tikzpicture}[scale=0.95, transform shape]
  % שכבה 1: מאקרו ורגולציה
  \fill[BrandBlue!20, draw=BrandNavy, line width=1pt]
    (-6.0, 4.2) -- (6.0, 4.2) -- (4.8, 3.15) -- (-4.8, 3.15) -- cycle;
  \node[font=\small\bfseries, text=BrandNavy] at (0, 3.68)
    {\texthebrew{1. מגמות מאקרו ורגולציה: זינוק ריבית (}\en{0.1\% $\to$ 4.75\%}\texthebrew{) והוראת השעה של בנק ישראל (}\en{2025}\texthebrew{)}};

  % שכבה 2: ההקשר היישומי
  \fill[BrandBlue!35, draw=BrandNavy, line width=1pt]
    (-4.7, 3.05) -- (4.7, 3.05) -- (3.6, 2.0) -- (-3.6, 2.0) -- cycle;
  \node[font=\small\bfseries, text=BrandNavy] at (0, 2.52)
    {\texthebrew{2. ההקשר היישומי: שוק הדיור בישראל --- }\en{>85,000}\texthebrew{ דירות לא-מכורות ומבצעי קבלן }\en{20/80}};

  % שכבה 3: המוטיבציה ופער הידע
  \fill[BrandGold!30, draw=BrandGold!90!black, line width=1pt]
    (-3.5, 1.9) -- (3.5, 1.9) -- (2.4, 0.85) -- (-2.4, 0.85) -- cycle;
  \node[font=\small\bfseries, text=BrandNavy] at (0, 1.38)
    {\texthebrew{3. המוטיבציה: היעדר מודל המאחד הלם תשלום (}\en{PTI}\texthebrew{) ומבחן קיצון ליזמים (}\en{DSCR}\texthebrew{)}};

  % שכבה 4: תוצר הדוח המכין
  \fill[BrandTeal!30, draw=BrandTeal, line width=1pt]
    (-2.3, 0.75) -- (2.3, 0.75) -- (1.5, -0.25) -- (-1.5, -0.25) -- cycle;
  \node[font=\small\bfseries, text=BrandNavy] at (0, 0.25)
    {\texthebrew{4. המחקר הנוכחי: סימולציית קיר המסירות }\en{2026--2027}};
\end{tikzpicture}
\end{english}
\caption{מבנה המחקר כמשפך: מהמגמה המאקרו-מימונית הרחבה אל ארכיטקטורת הסימולציה של הדוח המכין (בהשראת מתווה ד"ר סגל \cite{segal2026prep})}
\label{fig:ch01_funnel}
\end{figure}

\textbf{התחנה הראשונה במשפך (מגמות מאקרו-כלכליות ומוניטריות):} בין הרבעון הראשון של שנת \num{2022} לרבעון השלישי של שנת \num{2023} ביצעו הבנקים המרכזיים בעולם, ובראשם הבנק הפדרלי בארה"ב (\en{Federal Reserve}) ובנק ישראל (\en{BOI}), את מהלך ההידוק המוניטרי המהיר ביותר בארבעת העשורים האחרונים \cite{boi2024annual}. בישראל זינקה הריבית המוניטרית מרמת שפל היסטורית של \pct{0.1} במרץ \num{2022} לרמה של \pct{4.75} במאי \num{2023}, והתייצבה ברמה מרסנת של \pct{4.00} עד \pct{4.50} לאורך השנים \num{2024--2026} \cite{boi2025stability}. שינוי משטר הריבית ייקר באחת הן את עלות המשכנתאות למשקי הבית והן את עלות אשראי הליווי הבנקאי ליזמי הבנייה (המתומחר לרוב בריבית פריים בתוספת מרווח סיכון של \pct{1.0} עד \pct{2.5}).

\textbf{התחנה השנייה (ההקשר היישומי בשוק הדיור הישראלי):} בניגוד למחזורי האטה קודמים שבהם התכווצות הביקושים תורגמה לירידת מחירים נומינלית מהירה, שוק הדיור הראשוני בישראל בשנים \num{2022--2026} פיתח מנגנון עקיפה פיננסי: \textbf{מבצעי המימון הקבלניים}. יזמים הכבולים להסכמי ליווי בנקאי סגור (\en{Closed Project Accompaniment}) נמנעו מהפחתת "מחירי המדף" הנומינליים של הדירות, מחשש ששחיקת המחיר המוצהר תפגע ביחס ההלוואה לשווי (\en{LTV}) ובעודפי הפרויקט החזויים בדוח האפס. תחת זאת, העניקו היזמים לרוכשים הטבות מימון עמוקות: דחיית תשלומים במתכונת \num{20/80}, \num{15/85} או \num{10/90} (תשלום הון עצמי מינימלי בחתימה ודחיית היתרה למסירה ללא הצמדה למדד תשומות הבנייה), או העמדת הלוואות גישור ("בלון" / \en{Bullet}) שבהן היזם מסבסד את תשלומי הריבית לבנק המלווה במשך \num{36} עד \num{48} חודשי הבנייה \cite{boi2025stability,mof2025economist}.

\textbf{התחנה השלישית (המוטיבציה המחקרית):} בעקבות מדיניות זו צמח מלאי הדירות החדשות הלא-מכורות בידי קבלנים מרמה של \num{44,200} יחידות דיור בתחילת \num{2022} לשיא היסטורי של \num{87,400} יחידות דיור בתחילת \num{2026} (ו-\num{85,900} ברבעון השלישי של \num{2026}) \cite{cbs2026dwellings}. במקביל, עשרות אלפי משקי בית שרכשו דירות "על הנייר" בשנים \num{2023--2024} מתקרבים כעת למועד קבלת טופס \num{4} ומסירת החזקה בשנים \num{2026--2027} ("קיר המסירות"). במועד זה יידרשו הרוכשים להמיר את הלוואת הבלון הקבלנית או את יתרת החוב הדחויה (\pct{80} עד \pct{90} ממחיר הדירה) למשכנתא בנקאית קבועה בריבית שוק גבוהה, תחת מגבלות יחס החזר להכנסה (\en{PTI}) קשיחות של בנק ישראל \cite{boi2025directive}. הדיווחים הקיימים בוחנים את משקי הבית, היזמים והבנקים כגזרות נפרדות, אך חסרים מודל כמותי אינטגרטיבי הבוחן את שרשרת ההדבקה הדו-כיוונית ביניהם בהשוואה למשבר הסאב-פריים האמריקאי של שנת \num{2008} \cite{mian2009consequences,mayer2009rise}.

\textbf{התחנה הרביעית (מבנה הדוח ותקציב העמודים):} כדי להבטיח איזון אנליטי ומניעת הטיה מבנית, תוכנן הדוח המכין לפי "מפת \num{40} העמודים" של ד"ר סגל \cite{segal2026prep,segal2026l03}. טבלה~\ref{tab:ch01_page_budget} מפרטת את מתווה הקצאת העמודים בין החלקים המקדמיים, שבעת פרקי הליבה ורשימת המקורות.

\begin{table}[htbp]
\centering
\caption{מתווה הקצאת העמודים בדוח המכין (בהתאם למפת העמודים של ד"ר סגל \cite{segal2026prep})}
\label{tab:ch01_page_budget}
\small
\begin{tabularx}{\textwidth}{R{3.3cm} Y C{1.6cm} C{2.2cm}}
\toprule
\rowcolor{TableHeaderBg}
\textbf{פרק / רכיב במסמך} & \textbf{תכולה מרכזית ותוצרים אנליטיים} & \textbf{היקף יעד} & \textbf{מספור עמודים} \\
\midrule
חלקים מקדמיים & דף שער ממלכתי, זכויות, תקציר מנהלים (\num{250--350} מילים), מילות מפתח, תוכן עניינים, רשימות איורים וטבלאות, קיצורים. & \num{6} עמ' & \en{i--vi} \\
\addlinespace[2pt]
פרק \num{1}: מסגרת המחקר & מבנה משפך, הגדרת בעיה כמותית, שאלת מחקר והשערות (\en{$H_1$--$H_3$}), מפרט דרישות (\en{REQ-RES}), מטריצת ספרות ומאמר הדגל. & \num{6} עמ' & \num{1--6} \\
\addlinespace[2pt]
פרק \num{2}: ארה"ב \num{2008} & אנטומיה של משבר הסאב-פריים (\num{2001--2008}): הלוואות \en{ARM} ו-\en{NINJA}, שרשרת האיגוח (\en{MBS/CDO}) וקיר איפוס הריביות. & \num{6} עמ' & \num{7--12} \\
\addlinespace[2pt]
פרק \num{3}: ישראל \num{2022--2026} & שוק הדיור בישראל, מלאי \num{85,000+} דירות, אבולוציית מבצעי \num{20/80} ו-\num{20/50/30}, פילוח \en{TAM/SAM/SOM} ומודל הנחת ה-\en{NPV}. & \num{6} עמ' & \num{13--18} \\
\addlinespace[2pt]
פרק \num{4}: ניתוח השוואתי & השוואה רב-פרמטרית (ישראל \num{2026} מול ארה"ב \num{2008}), תרשים רדאר, מוקדי דמיון מול חוק המכר, \en{Recourse} ומגבלות \en{LTV/PTI}. & \num{6} עמ' & \num{19--24} \\
\addlinespace[2pt]
פרק \num{5}: מתודולוגיה ותכן & מאגרי נתונים (כרמ"ן, למ"ס, בנק ישראל), מטריצת החלטה (\en{Pugh}), משוואות \en{Payment Shock} ו-\en{Stress Testing} ב-\en{Python}. & \num{6} עמ' & \num{25--30} \\
\addlinespace[2pt]
פרק \num{6}: קיר המסירות & ממצאי סימולציית הקיצון: התפלגות \en{PTI} למשקי בית, שחיקת \en{DSCR} של יזמים ממונפים (\en{Mezzanine}) ועמידות הון הבנקים (\en{CET1}). & \num{5} עמ' & \num{31--35} \\
\addlinespace[2pt]
פרק \num{7}: תוכנית וסיכונים & מבנה פירוק עבודה (\en{WBS}), מטריצת \en{RACI}, תרשים גאנט (\num{14} שבועות), מטריצת סיכונים (\en{$R=P\times I$}) ומוכנות למימוש. & \num{4} עמ' & \num{36--39} \\
\addlinespace[2pt]
מקורות ביבליוגרפיים & רשימת מקורות שפיטים, דוחות מוסדיים ותקנים בפורמט \en{IEEE} המנוהלת אוטומטית באמצעות \en{biber}. & \num{2} עמ' & \num{40--41} \\
\bottomrule
\end{tabularx}
\end{table}

\section{ניסוח כמותי של הבעיה, שאלת המחקר והשערות המחקר}
\label{sec:ch01_problem}

על פי סעיף \num{4.1} במדריך של ד"ר סגל \cite{segal2026prep}, הגדרת בעיה מקצועית אינה נשענת על אמירות איכותניות כגון "שוק הנדל"ן מצוי בקיפאון מסוכן", אלא מציגה פער מדיד וחד-משמעי בין המצב המצוי לבין ספי היציבות הפיננסית. איור~\ref{fig:ch01_problem_chain} מתאר את שרשרת הגזירה הלוגית מהמצב המצוי אל מפרט הדרישות של המחקר.

\begin{figure}[htbp]
\centering
\begin{english}
\begin{tikzpicture}[
  node distance=0.65cm,
  box/.style={rectangle, rounded corners=4pt, draw=BrandNavy, line width=1pt, minimum width=2.65cm, minimum height=1.55cm, align=center, font=\footnotesize\bfseries}
]
  \node[box, fill=BrandGray] (curr) at (9.6, 0)
    {\texthebrew{המצב המצוי}\\ \en{85.9K}\ \texthebrew{לא-מכורות}\\ \en{38\%--46\%}\ \texthebrew{עסקאות בלון}};
  \node[box, fill=BrandLightGold] (gap) at (6.4, 0)
    {\texthebrew{הפער האנליטי}\\ \texthebrew{הנחת }\en{NPV}\ \texthebrew{סמויה}\\ \texthebrew{וחיתום דחוי למסירה}};
  \node[box, fill=BrandLightBlue] (rq) at (3.2, 0)
    {\texthebrew{שאלה והשערות}\\ \texthebrew{אמידת }\en{$\Delta\text{PTI}$}\texthebrew{, }\en{DSCR}\\ \texthebrew{ויציבות }\en{CET1}};
  \node[box, fill=BrandLightTeal] (req) at (0, 0)
    {\texthebrew{מפרט דרישות}\\ \en{REQ-RES-01..08}\\ \texthebrew{יעדים ושיטות אימות}};

  \draw[-Stealth, line width=1.2pt, color=BrandNavy] (curr) -- (gap);
  \draw[-Stealth, line width=1.2pt, color=BrandNavy] (gap) -- (rq);
  \draw[-Stealth, line width=1.2pt, color=BrandNavy] (rq) -- (req);
\end{tikzpicture}
\end{english}
\caption{שרשרת גזירת הבעיה: מהמצב המצוי בשוק הדיור והאשראי אל מפרט הדרישות הכמותי (לפי איור \num{4.1} בספר של ד"ר סגל \cite{segal2026prep})}
\label{fig:ch01_problem_chain}
\end{figure}

\subsection{ניסוח כמותי של מרחב הבעיה}
בחינת נתוני בנק ישראל, הלמ"ס ואגף הכלכלן הראשי לשנים \num{2022--2026} \cite{boi2025stability,cbs2026dwellings,mof2025economist} מעלה ארבע עובדות כמותיות המגדירות את הבעיה:
\begin{enumerate}[leftmargin=*, itemsep=3pt]
  \item \textbf{היקף חשיפת עסקאות הקבלן הדחויות:} בין מחצית \num{2023} למרץ \num{2025} הגיע שיעור העסקאות בשוק הדיור החדש שכללו הטבות מימון קבלניות (מבצעי \num{20/80}, \num{10/90} והלוואות בולט/בלון מסובסדות) ל-\pct{38} עד \pct{46} מכלל העסקאות הארציות, ולמעלה מ-\pct{55} באזורי הביקוש במרכז הארץ ובדרומה \cite{boi2025stability}. בסך הכול מוערכת מצבת החוזים הדחויים הצפויים להגיע לאכלוס בשנים \num{2026--2027} בכ-\num{31,500} משקי בית, בהיקף אשראי דחוי מצטבר של כ-\num{68} מיליארד ש"ח.
  \item \textbf{עיוות המחיר הנומינלי מול הערך הכלכלי המהוון (\en{NPV}):} מדד מחירי הדירות של הלמ"ס מודד את המחיר החוזי הנומינלי המדווח לרשות המסים, ואינו מנכה את שווי סבסוד הריבית ופטור ההצמדה למדד תשומות הבנייה שסופג היזם. בהיוון לפי עלות ההון האפקטיבית של היזמים (\pct{5.5} עד \pct{6.5} לשנה בתוספת ויתור על הצמדה של \pct{2.5} לשנה), מחיר החוזה הנומינלי מנופח ב-\pct{8.5} עד \pct{12.4} מעל הערך הכלכלי הריאלי של העסקה (\en{Implicit NPV Discount}).
  \item \textbf{הלם התשלום (\en{Payment Shock}) במועד המסירה:} רוכש טיפוסי שחתם על עסקת \num{20/80} בדירה שמחירה \num{2.65} מיליון ש"ח (יתרת תשלום של \num{2.12} מיליון ש"ח במסירה) תחת הנחה שהריבית תרד ל-\pct{3.4}, צפה החזר חודשי של כ-\nis{10,490}. בפועל, בהינתן ריבית משכנתא אפקטיבית של \pct{5.25} ותוספת עלויות הצמדה ושדרוגים של \pct{6.8}, מזנק ההחזר החודשי במועד המסירה ל-\nis{13,570} --- תוספת הלם תשלום ממוצעת של \nis{3,080} לחודש.
  \item \textbf{צוואר הבקבוק היזמי והבנקאי:} מלאי של \num{85,900} דירות לא-מכורות מחייב את היזמים לממן הוצאות ריבית שוטפות בגין אשראי ליווי וקרקע בהיקף חסר תקדים. חברות יזמיות מהשורה השנייה והשלישית, שהשלימו הון עצמי לפרויקטים באמצעות הלוואות מזנין (\en{Mezzanine}) חוץ-בנקאיות בריבית של \pct{10.5} עד \pct{13.5}, חשופות לשחיקת יחס כיסוי החוב (\en{DSCR}) מתחת לסף חדלות הפירעון (\num{1.00x}) כבר בתרחיש של \pct{10} עד \pct{20} עיכובים או ביטולי חוזים מצד רוכשים שיידחו בחיתום הבנקאי במסירה.
\end{enumerate}

\subsection{שאלת המחקר המרכזית, שאלות המשנה והשערות המחקר}
בהתאם להבחנה שבטבלה \num{4.1} במדריך של ד"ר סגל \cite{segal2026prep} בין פרויקט מחקרי (הבוחן קשר סיבתי בין פרמטרים באמצעות השערות כמותיות) לבין פרויקט תכן מערכתי (המוגדר לפי יעדי \en{SMART} \cite{doran1981smart}), מחקר זה משלב את שתי הגישות:

\begin{modelbox}[שאלת המחקר המרכזית]
\textbf{"באיזו מידה מבצעי המימון הקבלניים, הלוואות הבלון ומינוף היזמים בשוק הנדל"ן הישראלי (\num{2022--2026}) משחזרים את מנגנוני הכשל של משבר הסאב-פריים (\num{2008}), ומהי עמידות המערכת הבנקאית והיזמית בישראל בפני גל חדלות פירעון בעת 'קיר המסירות' של \num{2026--2027}?"}
\end{modelbox}

משאלת-על זו נגזרות שלוש **שאלות משנה מחקריות** ושלוש **השערות מחקר כמותיות (\en{Testable Hypotheses})** הניתנות לאישוש או להפרכה אמפירית:

\begin{itemize}[leftmargin=*, itemsep=4pt]
  \item \textbf{שאלת משנה 1 (עיוות התמחור והלם התשלום למשקי הבית):} מהו שיעור ההנחה הכלכלית הסמויה (\en{$\Delta V_{\text{NPV}}$}) המגולם במבצעי המימון הקבלניים (\num{20/80}, \num{10/90} ו-\num{20/50/30}), ואיזה שיעור ממשקי הבית שרכשו דירות במבצעים אלו יחווה במועד המסירה הלם תשלום (\en{Payment Shock}) שיביא לחציית תקרת יחס ההחזר להכנסה הרגולטורית (\en{$\text{PTI} > 40\%$})?
  \item \textbf{השערת מחקר ראשונה (\en{$H_1$} --- השערת הלם התשלום וחסימת החיתום):} מבצעי הבלון הקבלניים (\num{20/80} ו-\num{15/85}) מגלמים הנחת \en{NPV} סמויה של \pct{8.0} עד \pct{11.0} ממחיר המחירון הנומינלי. בעת "קיר המסירות" (\num{2026--2027}), לפחות \pct{50} מרוכשי הבלון יחוו עלייה של מעל \pct{25} בהחזר החודשי האפקטיבי לעומת תחזית החתימה, ולפחות \pct{18} מהם יחצו את תקרת החיתום הבנקאית הקשיחה (\en{$\text{PTI} \ge 50\%$}) ויידרשו למימוש הנכס, להזרמת הון חיצוני או להסדר חוב.

  \item \textbf{שאלת משנה 2 (פגיעות היזמים תחת מבחני קיצון):} מהי רגישות יחס כיסוי החוב (\en{DSCR}) וכרית ההון העצמי של חברות בנייה יזמיות בישראל (בפילוח לפי גודל החברה ומבנה המינוף הבנקאי והחוץ-בנקאי שלה) לתרחישי ביטול או הפרת חוזים מצד רוכשים בשיעורים של \pct{10}, \pct{20} ו-\pct{30} במועד המסירה?
  \item \textbf{השערת מחקר שנייה (\en{$H_2$} --- השערת הפיצול היזמי וכשל אשראי המזנין):} בעוד שחברות בנייה ציבוריות גדולות (\en{Tier 1}, בעלות יחס הון עצמי למאזן של מעל \pct{20}) ישמרו על יחס כיסוי חוב חיובי (\en{$\text{DSCR} > 1.00\text{x}$}) גם בתרחיש קיצון של \pct{25} ביטולי עסקאות וירידת מחירים של \pct{10}, יזמים קטנים ובינוניים הנשענים על השלמת הון חוץ-בנקאית (\en{Mezzanine} בריבית דו-ספרתית, \en{Tier 3}) ייקלעו לחדלות פירעון תזרימית (\en{$\text{DSCR} < 1.00\text{x}$}) כבר בשיעור ביטולים או דחיית מסירות של \pct{10} עד \pct{12}.

  \item \textbf{שאלת משנה 3 (עמידות המערכת הבנקאית ומנגנוני הבלימה המוסדיים):} באיזו מידה מנגנוני ההגנה המבניים בישראל --- חוק המכר (דירות), ליווי בנקאי סגור, הלוואות \en{Recourse} אישיות, היעדר איגוח סאב-פריים (\en{MBS/CDO}) וכריות ההון של הבנקים (\en{CET1}) --- ימנעו את הפיכתו של משבר נזילות יזמי למשבר בנקאי מערכתי כדוגמת ארה"ב \num{2008}?
  \item \textbf{השערת מחקר שלישית (\en{$H_3$} --- השערת החסינות המערכתית הבנקאית):} בניגוד לקריסת הבנקים בארה"ב ב-\num{2008}, גם בתרחיש המאקרו-קיצוני החמור ביותר (\pct{30} כשל בעסקאות הבלון בשילוב ירידה של \pct{20} במחירי הדירות), יחס ההון העצמי רובד \num{1} (\en{CET1}) המצרפי של חמשת הבנקים הגדולים בישראל לא ירד מתחת לרף הדרישה הפיקוחית המינימלית (\pct{9.0}), הודות לכרית הספיגה הכפולה של הון הפרויקט בחשבון הליווי הסגור והערבות האישית (\en{Full Recourse}) של משקי הבית.
\end{itemize}

\section{מפרט הדרישות של המחקר: החוזה ההנדסי-אנליטי (\en{REQ-RES})}
\label{sec:ch01_requirements}

על פי פרק \num{4.3} בספרו של ד"ר סגל \cite{segal2026prep}, "דרישה בלי שיטת בדיקה היא משאלה". כדי להבטיח שהדוח המכין ישמש חוזה ביצוע מדויק ומחייב, תורגמו יעדי המחקר למפרט דרישות כמותי (\en{Requirements Specification}) המבחין בין דרישות תפקודיות-אנליטיות (\en{Functional / Analytical Requirements}) לבין אילוצי תכן ומתודולוגיה (\en{Constraints}). טבלה~\ref{tab:ch01_requirements} מציגה את מפרט הדרישות המלא של הפרויקט.

\begin{table}[htbp]
\centering
\caption{מפרט הדרישות הכמותי של פרויקט המחקר: החוזה ההנדסי-אנליטי (לפי טבלה \num{4.2} במדריך של ד"ר סגל \cite{segal2026prep})}
\label{tab:ch01_requirements}
\footnotesize
\begin{tabularx}{\textwidth}{R{2.1cm} R{3.2cm} C{1.5cm} R{2.8cm} Y}
\toprule
\rowcolor{TableHeaderBg}
\textbf{מזהה דרישה} & \textbf{תיאור הדרישה} & \textbf{סיווג} & \textbf{ערך יעד כמותי (\en{SMART})} & \textbf{שיטת אימות ותיקוף} \\
\midrule
\en{REQ-RES-01} & כיסוי אמפירי של עסקאות דירות חדשות ומלאי לא-מכור בישראל & נתונים & מדגם של \num{18} רבעונים (\num{2022Q1--2026Q2}), \num{N \ge 45,000} עסקאות & הצלבת קובצי כרמ"ן (רשות המסים) מול סדרות הלמ"ס ובנק ישראל \cite{cbs2026dwellings,ita2025carman}. \\
\addlinespace[2pt]
\en{REQ-RES-02} & מודל היוון (\en{NPV}) לחילוץ ההנחה הכלכלית הסמויה במבצעי קבלן & אנליטית & אמידת \en{$\Delta V_{\text{NPV}}$} בדיוק של \pct{\pm 0.25} עבור \num{4} מתווים (\num{20/80}, \num{15/85}, \num{10/90}, \num{20/50/30}) & הרצת מודל היוון תזרימי ב-\en{Python} מול תנאי המימון בדוחות בנק ישראל \cite{boi2025stability}. \\
\addlinespace[2pt]
\en{REQ-RES-03} & סימולציית \en{Monte Carlo} להלם תשלום (\en{Payment Shock}) למשקי בית & תפקודית & \num{N = 10,000} משקי בית; שגיאת תקן של אומדן \en{PTI} קטנה מ-\pct{0.3} & בדיקת התכנסות מונטה-קרלו (\en{Seed=20261115}) והשוואה להתפלגות ההכנסות של הלמ"ס. \\
\addlinespace[2pt]
\en{REQ-RES-04} & מודל מבחני קיצון (\en{Stress Test}) ליחס כיסוי חוב (\en{DSCR}) של יזמים & תפקודית & בחינת \num{3} קבוצות יזמים (\en{Tiers 1--3}) תחת \num{4} שיעורי ביטולים (\pct{0, 10, 20, 30}) & אימות מול דוחות כספיים מבוקרים של \num{25} חברות בנייה בורסאיות ב-\en{TASE}. \\
\addlinespace[2pt]
\en{REQ-RES-05} & מודל ספיגת הפסדי אשראי ויחס הלימות הון בנקאי (\en{CET1}) & תפקודית & סימולציה רבעונית על פני \num{8} רבעונים (\en{T0} עד \en{T+8Q}) מול רף פיקוחי של \pct{9.0} & השוואה למתודולוגיית מבחני הקיצון של חטיבת היציבות הפיננסית בבנק ישראל \cite{boi2025stability}. \\
\addlinespace[2pt]
\en{REQ-RES-06} & מטריצת השוואה רב-פרמטרית: ישראל \num{2026} מול ארה"ב \num{2008} & מחקרית & השוואה כמותית על פני \num{10} פרמטרים מוסדיים, משפטיים ומימוניים & תיקוף מול נתוני ה-\en{Fed} ומאמר הדגל של \en{Mian \& Sufi (2009)} \cite{mian2009consequences}. \\
\addlinespace[2pt]
\en{REQ-RES-07} & איכות עימוד, כיווניות (\en{BiDi}) וקומפילציה דטרמיניסטית ב-\en{LuaLaTeX} & אילוץ תכן & היקף של \num{40 \pm 5} עמודים; \num{0} שגיאות קומפילציה; \num{0} חריגות שוליים מעל \en{5pt} & הרצת \num{4} סבבי \en{lualatex + biber} ובדיקת דוח הדטקטור (\en{Prompt B2}) \cite{segal2026l03}. \\
\addlinespace[2pt]
\en{REQ-RES-08} & לוח זמנים וניהול סיכונים כמותי למימוש המחקר & ניהולית & תרשים גאנט של \num{14} שבועות; כימות \en{$R = P \times I$} ותוכנית מגירה לכל \en{$R \ge 12$} & בקרת אבני דרך (\en{Milestones}) ואימות שבוע רזרבה לפני כל אבן דרך \cite{segal2026prep}. \\
\bottomrule
\end{tabularx}
\end{table}

\section{סקירת ספרות תמטית וביקורתית, מטריצת ספרות ומאמר הדגל}
\label{sec:ch01_literature}

בהתאם להנחיות פרק \num{5} בספרו של ד"ר סגל \cite{segal2026prep} ומאמר המתודולוגיה של \en{Webster \& Watson (2002)} \cite{webster2002analyzing}, סקירת הספרות בדוח זה אינה מסכמת מאמרים ברצף כרונולוגי, אלא מחלקת את הידע המדעי לארבעה צירי ליבה תמטיים, מנתחת ביקורתית את ההנחות המפשטות של כל זרם מחקרי, ומזהה את **פער הידע (\en{Research Gap})** באמצעות מטריצת ספרות ומאמר דגל.

\subsection{ציר תמטי א': אסימטריית מידע, בחירה שלילית וקיצוב אשראי}
הבסיס התיאורטי לניתוח כשלים בשוקי האשראי לדיור הונח במודל הקלאסי של \en{Stiglitz \& Weiss (1981)} \cite{stiglitz1981credit}. המודל מוכיח כי בשוק אשראי המאופיין באסימטריית מידע (\en{Information Asymmetry}) בין המלווה ללווה, העלאת שער הריבית מעבר לנקודת שיווי משקל קריטית אינה מגדילה את תוחלת התשואה של הבנק, משום שהיא מפעילה שני מנגנונים הרסניים:
\begin{enumerate}[leftmargin=*, itemsep=2pt]
  \item \textbf{בחירה שלילית (\en{Adverse Selection}):} לווים סולידיים בעלי הסתברות החזר גבוהה יוצאים מהשוק כאשר הריבית מתייקרת, בעוד שלווים ספקולטיביים בעלי תיאבון סיכון גבוה מוכנים לחתום על חוזי הלוואה בריבית גבוהה או בתנאי דחיית תשלום קיצוניים.
  \item \textbf{סיכון מוסרי (\en{Moral Hazard}) והסטת סיכונים (\en{Risk Shifting}):} כפי שהראו \en{Jensen \& Meckling (1976)} \cite{jensen1976theory}, כאשר שיעור ההון העצמי של הלווה או היזם בעסקה נמוך מאוד (מינוף גבוה), פונקציית התשואה שלו הופכת לאופציית רכש (\en{Call Option}) על ערך הנכס: אם מחיר הנכס עולה, הלווה גורף את מלוא רווח ההון הממונף; אם המחיר יורד, ההפסד מוגבל להון העצמי הדל או מתגלגל אל המלווה.
\end{enumerate}
בישראל של השנים \num{2022--2026}, מבצעי ה-\num{10/90} וה-\num{20/80} יצרו בדיוק את עיוות התמריצים הזה: רוכשים בעלי הון עצמי נמוך או משקיעים ספקולטיביים ("פליפרים") נכנסו לעסקאות רכישה של דירות בשווי \num{2.5} עד \num{4.0} מיליון ש"ח תוך סיכון הון ראשוני של \num{250} עד \num{400} אלף ש"ח בלבד, בהנחה שיוכלו למכור את הזכות החוזית לפני המסירה או למחזר את החוב בריבית נמוכה יותר \cite{boi2025stability,mof2025economist}.

\subsection{ציר תמטי ב': השערת אי-היציבות הפיננסית ומחזורי אשראי בנדל"ן}
היימן מינסקי (\en{Minsky, 1992}) \cite{minsky1992financial} פיתח את **השערת אי-היציבות הפיננסית (\en{Financial Instability Hypothesis})**, המבחינה בין שלושה משטרי מימון המתפתחים לאורך מחזור גאות בנדל"ן:
\begin{enumerate}[leftmargin=*, itemsep=2pt]
  \item \textbf{מימון מגודר (\en{Hedge Financing}):} תזרימי ההכנסה השוטפים של הלווה מכסים במלואם הן את תשלומי הריבית והן את פירעון קרן החוב בכל תקופה.
  \item \textbf{מימון ספקולטיבי (\en{Speculative Financing}):} תזרימי ההכנסה מכסים את הריבית השוטפת בלבד, אך הלווה נדרש לגלגל (\en{Roll Over}) את קרן החוב במועד הפירעון.
  \item \textbf{מימון פונזי (\en{Ponzi Financing}):} תזרימי ההכנסה השוטפים אינם מכסים אפילו את תשלומי הריבית או הקרן, והיכולת לעמוד בהתחייבויות נשענת בלעדית על המשך עליית מחירי הנכסים (\en{Asset Price Appreciation}) או על הזרמת אשראי חדש ברגע המיחזור.
\end{enumerate}
מחקריהם המקיפים של \en{Kindleberger \& Aliber (2011)} \cite{kindleberger2011manias} ושל \en{Reinhart \& Rogoff (2009)} \cite{reinhart2009this} מוכיחים כי משברי נדל"ן ובנקאות היסטוריים פורצים בדיוק בנקודת המעבר ממשטר מימון ספקולטיבי למשטר מימון פונזי ("רגע מינסקי" --- \en{Minsky Moment}), כאשר זעזוע ריבית מוניטרי עוצר את עליית מחירי הנכסים וחוסם את ערוץ המחזור במועד פקיעת הלוואות הגישור.

\subsection{ציר תמטי ג': איגוח, שחיקת החיתום וניפוח מחירים באמצעות הטבות מוכר}
בבחינת משבר הסאב-פריים האמריקאי של \num{2008}, הראו \en{Keys et al.\ (2010)} \cite{keys2010did} ו-\en{Gorton (2009)} \cite{gorton2009securitized} כי מעבר הבנקים ממודל של החזקת הלוואות במאזן (\en{Originate-to-Hold}) למודל של הפצה ואיגוח (\en{Originate-to-Distribute}) באמצעות מכשירי \en{MBS} ו-\en{CDO} שחק את תמריצי הבנקים לבצע חיתום קפדני. במקביל, מחקרו פורץ הדרך של פרופ' יצחק בן-דוד (\en{Ben-David, 2011}) \cite{ben2011inflated} בכתב העת \en{American Economic Journal} הדגים מנגנון קריטי הזהה במהותו למבצעי הקבלנים בישראל: כאשר מוכרי דירות ויזמים מעניקים לרוכשים תמריצי מימון או החזרי מזומן סמויים (\en{Seller-Funded Concessions}) כדי לאפשר לרוכשים לעקוף מגבלות הון עצמי, מחיר העסקה הנומינלי המדווח מתנפח באופן מלאכותי ב-\pct{5} עד \pct{10} מעל שווי השוק הכלכלי האמיתי של הדירה, מה שמייצר אשליית יציבות במדדי המחירים הרשמיים ומגדיל את שיעור הכשל של הלווים בהמשך.

\subsection{מטריצת הספרות (\en{Literature Matrix}) ובחירת מאמר הדגל}
כדי למפות באופן שיטתי את תרומת הספרות הקיימת ולבודד את פער הידע (סעיף \num{5.1} במדריך של ד"ר סגל \cite{segal2026prep}), מרכזת טבלה~\ref{tab:ch01_lit_matrix} את ששת המקורות המרכזיים מול חמשת ממדי הניתוח של הפרויקט.

\begin{table}[htbp]
\centering
\caption{מטריצת ספרות השוואתית וזיהוי פער הידע המחקרי (לפי מתווה טבלה \num{5.1} בספר של ד"ר סגל \cite{segal2026prep})}
\label{tab:ch01_lit_matrix}
\footnotesize
\begin{tabularx}{\textwidth}{R{2.9cm} C{2.0cm} C{2.0cm} C{2.1cm} C{2.1cm} Y}
\toprule
\rowcolor{TableHeaderBg}
\textbf{מקור / מחקר} & \textbf{התרחבות אשראי והלם תשלום (\en{PTI})} & \textbf{ניפוח מחיר נומינלי מול \en{NPV}} & \textbf{מבחן קיצון ליזמים ואשראי מזנין} & \textbf{הגנות מוסדיות (חוק המכר ו-\en{Recourse})} & \textbf{סימולציה משולבת לקיר המסירות \num{2026}} \\
\midrule
\en{Mian \& Sufi (2009)} \cite{mian2009consequences} \textbf{(מאמר הדגל)} & \textbf{נבחן לעומק} & חלקי & לא נבחן & לא נבחן (\en{Non-Recourse}) & לא נבחן \\
\addlinespace[2pt]
\en{Mayer et al.\ (2009)} \cite{mayer2009rise} & \textbf{נבחן לעומק} & חלקי & לא נבחן & חלקי & לא נבחן \\
\addlinespace[2pt]
\en{Ben-David (2011)} \cite{ben2011inflated} & חלקי & \textbf{נבחן לעומק} & לא נבחן & לא נבחן & לא נבחן \\
\addlinespace[2pt]
\en{Ghent \& Kudlyak (2011)} \cite{ghent2011recourse} & חלקי & לא נבחן & לא נבחן & \textbf{נבחן לעומק} (\en{Recourse}) & לא נבחן \\
\addlinespace[2pt]
דוח יציבות בנק ישראל (\num{2025}) \cite{boi2025stability} & חלקי (מאקרו) & חלקי & חלקי (בנקאי בלבד) & \textbf{נבחן לעומק} & חלקי (ללא מיקרו-סימולציה) \\
\addlinespace[2pt]
\rowcolor{BrandLightTeal}
\textbf{המחקר הנוכחי (\num{2026})} & \textbf{נבחן (\en{Monte Carlo})} & \textbf{נבחן (\en{NPV})} & \textbf{נבחן (\en{DSCR + Mezzanine})} & \textbf{נבחן (ישראל מול ארה"ב)} & \textbf{זה הפער (משולב ומכומת)} \\
\bottomrule
\end{tabularx}
\end{table}

\textbf{ניתוח מאמר הדגל (\en{Flagship / Benchmark Paper}):} על פי שלושת הקריטריונים של ד"ר סגל (סעיף \num{5.2} ואיור \num{5.1} במדריך \cite{segal2026prep}: פרסום בבמה שפיטה מובילה, קרבה ישירה לבעיה, וחזרתיות כמותית), נבחר מחקרם של \textbf{\en{Atif Mian \& Amir Sufi (2009)}}, \en{``The Consequences of Mortgage Credit Expansion: Evidence from the U.S.\ Mortgage Default Crisis''}, שפורסם בכתב העת המוביל \en{The Quarterly Journal of Economics} \cite{mian2009consequences}, כמאמר הדגל של הפרויקט (בשילוב המשוואה המבנית של \en{Ben-David, 2011} \cite{ben2011inflated}).

\en{Mian \& Sufi} ניתחו נתוני מיקרו של כ-\num{3,000} אזורי מיקוד (\en{ZIP Codes}) בארה"ב בין השנים \num{1996--2007}, והוכיחו באמצעות מודל הפרשים-בהפרשים (\en{Differences-in-Differences}) ומשתני עזר כי התרחבות היצע האשראי ללווים שוליים בעלי יחס החזר להכנסה גבוה הובילה תחילה לעלייה חדה במחירי הדירות הנומינליים, אך ברגע שהריבית התאפסה ב-\num{2007--2008} קרסו שיעורי הפירעון באותם אזורים בדיוק. המשוואה האקונומטרית השולטת במאמר הדגל היא:
\begin{equation}
\label{eq:mian_sufi_flagship}
\Delta \text{Default}_{z, t_0 \to t_1} = \alpha + \beta_1 \cdot \Delta \text{CreditExpansion}_{z, t_{-1} \to t_0} + \beta_2 \cdot \Delta \text{HousePrice}_{z} + \boldsymbol{\gamma}' \mathbf{X}_{z} + \varepsilon_{z}
\end{equation}
כאשר \en{$\Delta \text{Default}_{z, t_0 \to t_1}$} הוא השינוי בשיעור חדלות הפירעון באזור \en{$z$} בעת קיר איפוס הריביות, \en{$\Delta \text{CreditExpansion}$} הוא שיעור חדירת הלוואות המינוף והבלון בתקופת הגאות, ו-\en{$\mathbf{X}_{z}$} הוא וקטור בקרות דמוגרפיות ותעסוקתיות. עם זאת, המודל של \en{Mian \& Sufi} מניח סביבה אמריקאית של הלוואות \en{Non-Recourse}, איגוח מלא (\en{MBS}) וחיתום המבוצע במועד החתימה בלבד, ואינו מתחשב במבנה הישראלי הייחודי שבו **היזם עצמו** מממן את הלוואת הבלון בתקופת הבנייה תחת ליווי בנקאי סגור וחוק המכר \cite{salelaw1974}. מתוך ניתוח זה נגזר **פער הידע (\en{Research Gap})** של מחקרנו \cite{sandberg2011ways}:

\begin{modelbox}[פער הידע המחקרי (\en{Research Gap})]
\textbf{עד כה לא נבדק במודל סימולציה כמותי מאוחד כיצד אינטראקציה בין דחיית חיתום משקי הבית במבצעי קבלן (\num{20/80}) לבין מינוף חוץ-בנקאי (\en{Mezzanine}) של יזמים מתורגמת לשיעורי ביטול עסקאות, לשחיקת יחס כיסוי החוב היזמי (\en{DSCR}) ולפגיעה בהלימות ההון הבנקאית (\en{CET1}) תחת משטר הגנת חוק המכר והלוואות \en{Recourse} בישראל בעת "קיר המסירות" של \num{2026--2027}.}
\end{modelbox}

\section{מטרות המחקר, גבולות גזרה (\en{Scope}) ומדדי הצלחה}
\label{sec:ch01_scope}

\textbf{מטרות המחקר:} (א) לבנות מודל תמחור מהוון (\en{NPV}) המבודד את ההנחה הכלכלית האמיתית המגולמת במחירי הדירות החדשות בישראל בשנים \num{2022--2026}; (ב) לאמוד באמצעות סימולציית \en{Monte Carlo} את התפלגות הלם התשלום (\en{Payment Shock}) של משקי הבית במועד האכלוס; (ג) לבצע מבחני קיצון (\en{Stress Testing}) למאזני חברות הבנייה ולמערכת הבנקאית תחת תרחישי ביטול חוזים וירידות מחירים; (ד) לגבש המלצות מדיניות מבוססות-נתונים עבור הפיקוח על הבנקים, רשות המסים ורשות ניירות ערך.

\textbf{גבולות הגזרה (\en{Project Scope}):} המחקר מתמקד בשוק הדיור למגורים בישראל (בדגש על דירות חדשות מקבלנים שנמכרו בין \num{2022Q1} ל-\num{2026Q3} ועתידות להימסר בשנים \num{2026--2027}), בחמשת הבנקים הגדולים בישראל ובמדגם מייצג של חברות בנייה יזמיות (ציבוריות ופרטיות), אל מול אמת המידה ההיסטורית של משבר הסאב-פריים בארה"ב (\num{2001--2008}). המחקר אינו עוסק בנדל"ן מניב משרדי או מסחרי.

\begin{warningbox}[מלכודת מתודולוגית: בלבול בין מחיר חוזי נומינלי לבין שווי כלכלי ריאלי]
הסתמכות נאיבית על מדד מחירי הדירות הנומינלי בתקופה של מבצעי מימון אגרסיביים יוצרת הטיה כלפי מעלה באומדן שווי הבטוחות של הבנקים: כאשר קבלן מוכר דירה ב-\num{2.80} מיליון ש"ח במבצע \num{10/90} ללא הצמדה למדד ובסבסוד הלוואת קבלן ל-\num{3} שנים, השווי הכלכלי הנוכחי של העסקה (\en{NPV}) הוא כ-\num{2.48} מיליון ש"ח בלבד. כל מודל סיכון המתעלם מפער זה שוגה באמידת יחס ה-\en{LTV} האמיתי של הפרויקט.
\end{warningbox}
